\documentclass[11pt,a4paper]{article}

\usepackage[utf8]{inputenc}
\usepackage[T1]{fontenc}
\usepackage{amsmath,amssymb,amsthm}
\usepackage{graphicx}
\usepackage[margin=1in]{geometry}

\newcommand{\toprule}{\hline}
\newcommand{\midrule}{\hline}
\newcommand{\bottomrule}{\hline}
\newcommand{\url}[1]{\texttt{#1}}

\newtheorem{definition}{Definition}
\newtheorem{theorem}{Theorem}
\newtheorem{lemma}{Lemma}
\newtheorem{claim}{Claim}

\newcommand{\rotl}{\mathrm{rotl}_{64}}
\newcommand{\ctz}{\mathrm{ctz}}

\title{%
    \textbf{Collatz-Based Cryptographic Primitives:}\\[4pt]
    \large A Permutation, Hash Function, Authenticated Cipher,\\
    and Message Authentication Code\\
    Derived from the $3n+1$ Mapping
}

\author{
    John Mock\\[4pt]
    Zachary Mock\\[6pt]
    \small Draft for Peer Review --- March 2026
}

\date{}

\begin{document}
\maketitle

\begin{abstract}
We present a family of cryptographic primitives whose core nonlinear
operation is derived from the Collatz mapping
$n \mapsto (3n+1)/2^{\nu_2(3n+1)}$.
The central construction is \textbf{CollatzPerm-512}, a 512-bit
cryptographic permutation that pairs a branchless Collatz step with a
$64 \times 64 \to 128$-bit integer multiplication for diffusion. On top
of this single permutation we build three modes of operation using
well-studied sponge and duplex constructions:
(1)~\textbf{CollatzHash}---a sponge hash with 256-, 512-, or 1024-bit
output;
(2)~\textbf{CollatzAEAD}---a duplex-sponge authenticated cipher
(256-bit key, 128-bit nonce, 128-bit tag);
(3)~\textbf{CollatzMAC}---a keyed-sponge message authentication code.
The design targets 128-bit security. We provide a reference C++17
implementation with ARM64 inline assembly for the critical path.
Empirical results on Apple M3 Max show CollatzHash-256 at 17\,M\,ops/s
(3.4$\times$ faster than standalone SHA-256), CollatzAEAD at 860\,MB/s on
4\,KB messages, and CollatzMAC at 6.9\,M\,ops/s (5.4$\times$ faster than
HMAC-SHA-256). All 70 security validation tests pass, including full
permutation diffusion, avalanche at 49.6--50.3\%, zero collisions in
$10^6$ inputs, AEAD authentication integrity, and constant-time tag
comparison.
\end{abstract}

\tableofcontents
\newpage

% =====================================================================
\section{Introduction}
% =====================================================================

The Collatz conjecture (also known as the $3n+1$ problem) states that
for any positive integer $n$, the sequence defined by
\[
    T(n) = \begin{cases}
        n/2 & \text{if } n \equiv 0 \pmod{2},\\
        (3n+1)/2 & \text{if } n \equiv 1 \pmod{2},
    \end{cases}
\]
eventually reaches 1. Despite its elementary statement, the conjecture
remains one of the most famous open problems in mathematics. The mapping
exhibits chaotic, pseudorandom behavior: nearby starting values can
produce vastly different orbit lengths and trajectories.

We exploit this pseudorandom behavior as a \emph{nonlinear mixing
primitive} for cryptographic hash functions and ciphers. Prior work has
explored Collatz-like maps in pseudorandom number generation
\cite{lagarias1985}, but to our knowledge no complete cryptographic suite
has been built on the $3n+1$ mapping.

\subsection{Contributions}

\begin{enumerate}
    \item A novel cryptographic permutation, \textbf{CollatzPerm-512},
          that uses a branchless Collatz step paired with 128-bit
          integer multiplication as its nonlinear layer.
    \item Three modes of operation---hash, AEAD, MAC---built using
          standard sponge/duplex constructions, inheriting formal
          security bounds from the permutation.
    \item A reference implementation in C++17 with hand-tuned ARM64
          assembly for the critical path.
    \item Empirical security evaluation: diffusion, avalanche,
          distribution, collision search, and authentication integrity.
    \item Performance benchmarks against SHA-256, ChaCha20, and
          HMAC-SHA-256.
\end{enumerate}

\subsection{Design Philosophy}

The fastest non-cryptographic hashes (rapidhash \cite{rapidhash},
wyhash \cite{wyhash}, xxHash3 \cite{xxhash}) share a key insight: a
single $64 \times 64 \to 128$-bit multiply followed by XOR of the high
and low halves provides massive diffusion in 3--4 CPU cycles. On the
other hand, established cryptographic hashes (SHA-256, SHA-3) derive
their security from carefully designed round functions with proven
diffusion properties, at the cost of higher latency.

Our design occupies a middle ground: we replace the purely algebraic
multiply with a \emph{Collatz--multiply hybrid} that adds mathematical
nonlinearity from number theory at minimal cycle cost, then use the
multiply for diffusion. The result is faster than SHA-256 while providing
a cryptographic-strength construction.

\subsection{Notation}

\begin{itemize}
    \item $\nu_2(n)$: the 2-adic valuation of $n$ (number of trailing
          zeros in the binary representation).
    \item $\rotl(x, r)$: 64-bit left rotation by $r$ positions.
    \item $\ctz(n) = \nu_2(n)$: count trailing zeros.
    \item $\oplus$: bitwise exclusive-or (XOR).
    \item $\cdot_{128}$: $64 \times 64$-bit unsigned multiplication
          producing a 128-bit result, split into low ($\mathrm{lo}$) and
          high ($\mathrm{hi}$) 64-bit words.
    \item $s[0..7]$: the 512-bit state as eight 64-bit words.
    \item $R_a = 12$: number of rounds for initialization and finalization.
    \item $R_b = 8$: number of rounds for per-block processing.
\end{itemize}

% =====================================================================
\section{Prior Work: CollatzHash V4 (Non-Cryptographic)}
\label{sec:v4}
% =====================================================================

Before developing the cryptographic suite, we designed
\textbf{CollatzHash V4}, a high-performance non-cryptographic hash
optimized for maximum throughput.

\subsection{Core Primitive: \texttt{collatz\_rapid\_mix}}

The V4 mixing function operates on two 64-bit inputs $(a, b)$ and
produces a 64-bit output. The algorithm proceeds as follows:

\begin{quote}\tt
\begin{tabbing}
\hspace{2em}\=\hspace{2em}\=\kill
\textrm{\textbf{Function}} collatz\_rapid\_mix$(a, b)$:\\
\> 1. $a \leftarrow a \,|\, 1$ \textrm{\small (ensure odd, branchless)}\\
\> 2. $t \leftarrow 3a + 1$\\
\> 3. $t \leftarrow t \gg \ctz(t)$ \textrm{\small (first Collatz iteration)}\\
\> 4. $t_2 \leftarrow 3t + 1$\\
\> 5. $t_2 \leftarrow t_2 \gg \ctz(t_2)$ \textrm{\small (second Collatz iteration)}\\
\> 6. $a \leftarrow t_2$\\
\> 7. $(\mathrm{lo}, \mathrm{hi}) \leftarrow a \cdot_{128} b$\\
\> 8. $a \leftarrow a \oplus \mathrm{lo}$;\quad $b \leftarrow b \oplus \mathrm{hi}$\\
\> 9. \textbf{return} $a \oplus b$
\end{tabbing}
\end{quote}

\subsection{Architecture}

V4 uses a 192-bit state (three 64-bit accumulators $s_0, s_1, s_2$)
with:
\begin{itemize}
    \item \textbf{Short-input fast path} ($\leq 48$ bytes): 2--3 mix
          calls, no loops.
    \item \textbf{Large-input path}: 3 parallel accumulators processing
          48-byte blocks.
    \item \textbf{Sponge squeeze}: for output widths beyond 64 bits (up
          to 1024 bits), re-permuting the state and emitting 64 bits per
          squeeze.
\end{itemize}

\subsection{ARM64 Assembly Optimization}

On Apple M3 Max (ARM64), the V4 mixing function is implemented in 11--16
instructions using:
\begin{itemize}
    \item \texttt{ADD}+\texttt{LSL} for $3n$ (fused add-shift),
    \item \texttt{RBIT}+\texttt{CLZ} for $\ctz$ (bit-reverse then count
          leading zeros),
    \item \texttt{MUL}+\texttt{UMULH} for 128-bit multiply (low and
          high halves).
\end{itemize}

A single-iteration variant (V4b) uses 11 instructions and achieves
16\,M\,ops/s at 1024-bit output and 253\,M\,ops/s at 64-bit output.

\subsection{Limitations}

V4's 192-bit state provides only $2^{64}$ collision resistance
(capacity $= 128$ bits), making it unsuitable for cryptographic use. It
has no key schedule, no formal security claims, and no resistance to
algebraic or differential attacks. This motivated the design of the
cryptographic suite described in the following sections.

% =====================================================================
\section{CollatzPerm-512: The Cryptographic Permutation}
\label{sec:perm}
% =====================================================================

\subsection{State}

The permutation operates on a 512-bit state represented as eight 64-bit
words: $s = (s[0], s[1], \ldots, s[7])$.

The state is divided into:
\begin{itemize}
    \item \textbf{Rate} ($r = 256$ bits): words $s[0..3]$, used for
          absorbing input and squeezing output.
    \item \textbf{Capacity} ($c = 256$ bits): words $s[4..7]$, never
          directly exposed.
\end{itemize}

\subsection{Round Function}

Each round $i$ ($0 \leq i < R$) applies three layers in sequence.

\subsubsection{Layer 1: Nonlinear (Collatz--Multiply Pairs)}

Four independent pairs $(s[0],s[1])$, $(s[2],s[3])$, $(s[4],s[5])$,
$(s[6],s[7])$ are processed by the \texttt{collatz\_mul\_pair} operation:

\begin{quote}\tt
\begin{tabbing}
\hspace{2em}\=\hspace{2em}\=\kill
\textrm{\textbf{Function}} collatz\_mul\_pair$(a, b)$:\\
\> 1. $a \leftarrow a \oplus \rotl(a, 23)$ \textrm{\small (pre-mix rotation)}\\
\> 2. $a \leftarrow a \,|\, 1$ \textrm{\small (ensure odd)}\\
\> 3. $t \leftarrow 3a + 1$\\
\> 4. $t \leftarrow t \gg \ctz(t)$ \textrm{\small (Collatz step)}\\
\> 5. $a \leftarrow t$\\
\> 6. $(\mathrm{lo}, \mathrm{hi}) \leftarrow a \cdot_{128} b$ \textrm{\small (128-bit multiply)}\\
\> 7. $a \leftarrow a \oplus \mathrm{lo}$;\quad $b \leftarrow b \oplus \mathrm{hi}$ \textrm{\small (XOR feedback)}
\end{tabbing}
\end{quote}

The pre-mix rotation (step 1) ensures that bit 0 of the input---which
would otherwise be destroyed by the $a \,|\, 1$ operation---influences
the output through bit position 23. This was identified through
diffusion testing (Section~\ref{sec:diffusion}).

\subsubsection{Layer 2: Linear Diffusion}

Cross-pair mixing ensures information propagates between all four pairs
within approximately 3 rounds:

\begin{quote}\tt
\begin{tabbing}
\hspace{2em}\=\hspace{2em}\=\kill
\textrm{\textbf{Linear diffusion layer:}}\\
\> 1. $s[0] \leftarrow s[0] \oplus s[2]$;\quad $s[4] \leftarrow s[4] \oplus s[6]$\\
\> 2. $s[1] \leftarrow s[1] \oplus s[3]$;\quad $s[5] \leftarrow s[5] \oplus s[7]$\\
\> 3. $(s[0], s[2], s[4], s[6]) \leftarrow (s[2], s[4], s[6], s[0])$ \textrm{\small (word rotation)}\\
\> 4. $s[1] \leftarrow s[1] \oplus \rotl(s[5], 19)$\\
\> 5. $s[3] \leftarrow s[3] \oplus \rotl(s[7], 31)$\\
\> 6. $s[5] \leftarrow s[5] \oplus \rotl(s[1], 7)$\\
\> 7. $s[7] \leftarrow s[7] \oplus \rotl(s[3], 53)$
\end{tabbing}
\end{quote}

The rotation constants $(19, 31, 7, 53)$ are chosen to be pairwise
coprime to 64, avoiding aligned mixing patterns.

\subsubsection{Layer 3: Round Constant Injection}

$$s[0] \leftarrow s[0] \oplus \mathrm{RC}[i]$$

The twelve round constants $\mathrm{RC}[0..11]$ are derived from the
binary expansions of $\pi$, $e$, $\sqrt{2}$, and $\sqrt{3}$
(nothing-up-my-sleeve numbers). When the round count exceeds 12,
constants are reused cyclically: $\mathrm{RC}[i \bmod 12]$.

The round constants in hexadecimal are:
{\small
\begin{align*}
\mathrm{RC}[0]  &= \mathtt{0x243f6a8885a308d3} & \mathrm{RC}[1]  &= \mathtt{0xb7e151628aed2a6a} \\
\mathrm{RC}[2]  &= \mathtt{0x6a09e667f3bcc908} & \mathrm{RC}[3]  &= \mathtt{0xbb67ae8584caa73b} \\
\mathrm{RC}[4]  &= \mathtt{0x3c6ef372fe94f82b} & \mathrm{RC}[5]  &= \mathtt{0xa54ff53a5f1d36f1} \\
\mathrm{RC}[6]  &= \mathtt{0x510e527fade682d1} & \mathrm{RC}[7]  &= \mathtt{0x9b05688c2b3e6c1f} \\
\mathrm{RC}[8]  &= \mathtt{0x1f83d9abfb41bd6b} & \mathrm{RC}[9]  &= \mathtt{0x5be0cd19137e2179} \\
\mathrm{RC}[10] &= \mathtt{0x428a2f98d728ae22} & \mathrm{RC}[11] &= \mathtt{0x7137449123ef65cd}
\end{align*}
}

\subsection{Round Counts}

Following the dual-round approach of Ascon \cite{ascon}:
\begin{itemize}
    \item $R_a = 12$ rounds: initialization and finalization (full
          security margin).
    \item $R_b = 8$ rounds: per-block absorption (reduced, compensated
          by wide state).
\end{itemize}

\subsection{Security Claim}

\begin{claim}
Assuming CollatzPerm-512 behaves as a pseudorandom permutation on 512
bits, the sponge and duplex constructions provide:
\begin{itemize}
    \item Collision resistance: $2^{c/2} = 2^{128}$.
    \item Preimage resistance:
          $\min(2^c, 2^{\mathrm{output\_bits}}) = 2^{256}$ (for 256-bit
          output).
    \item AEAD confidentiality and integrity: up to
          $\min(2^{c/2}, 2^k) = 2^{128}$ queries.
\end{itemize}
\end{claim}

These bounds follow directly from the generic sponge security
proofs \cite{sponge_security} with capacity $c = 256$.

% =====================================================================
\section{CollatzHash V5: Sponge Hash}
\label{sec:hash}
% =====================================================================

\subsection{Construction}

CollatzHash V5 is a sponge hash \cite{keccak_sponge} over CollatzPerm-512
with rate $r = 256$ bits and capacity $c = 256$ bits.

\begin{quote}\tt
\begin{tabbing}
\hspace{2em}\=\hspace{2em}\=\kill
\textrm{\textbf{Function}} CollatzHash$(M, d)$:
\textrm{\small Hash message $M$ to $d$ output bits}\\[4pt]
\textrm{\textit{Initialization:}}\\
\> 1. $s \leftarrow (0, 0, 0, 0, 0, 0, 0, 0)$ \textrm{\small (512-bit zero state)}\\
\> 2. $s[7] \leftarrow s[7] \oplus \mathtt{DS\_HASH}$ \textrm{\small (domain separation)}\\
\> 3. $s[6] \leftarrow s[6] \oplus d$ \textrm{\small (encode output length)}\\[4pt]
\textrm{\textit{Absorption:}}\\
\> 4. Pad $M$ with $10^*1$ padding to a multiple of $r = 256$ bits\\
\> 5. \textbf{for each} $r$-bit block $B_i$ of padded $M$:\\
\> \> $s[0..3] \leftarrow s[0..3] \oplus B_i$\\
\> \> $s \leftarrow \mathrm{CollatzPerm}(s, R_b)$\\[4pt]
\textrm{\textit{Finalization and Squeeze:}}\\
\> 6. $s \leftarrow \mathrm{CollatzPerm}(s, R_a)$ \textrm{\small (finalization)}\\
\> 7. $Z \leftarrow \varepsilon$ \textrm{\small (empty output)}\\
\> 8. \textbf{while} $|Z| < d$:\\
\> \> $Z \leftarrow Z \,\|\, s[0..3]$ \textrm{\small (squeeze $r$ bits)}\\
\> \> \textbf{if} $|Z| < d$ \textbf{then}
$s \leftarrow \mathrm{CollatzPerm}(s, R_b)$\\
\> 9. \textbf{return} first $d$ bits of $Z$
\end{tabbing}
\end{quote}

\subsection{Output Widths}

\begin{itemize}
    \item \textbf{256-bit}: 1 squeeze (no additional permutation after
          finalization).
    \item \textbf{512-bit}: 2 squeezes (1 additional permutation).
    \item \textbf{1024-bit}: 4 squeezes (3 additional permutations).
\end{itemize}

Each additional 256-bit block is a genuine squeeze from a re-permuted
state, following the same principle as SHA-3/SHAKE \cite{sha3}.

% =====================================================================
\section{CollatzAEAD: Authenticated Encryption}
\label{sec:aead}
% =====================================================================

CollatzAEAD is a duplex-sponge AEAD cipher \cite{duplex} with:
\begin{itemize}
    \item Key: 256 bits.
    \item Nonce: 128 bits.
    \item Tag: 128 bits.
    \item Rate: 256 bits per block.
\end{itemize}

\subsection{Construction}

\begin{quote}\tt
\begin{tabbing}
\hspace{2em}\=\hspace{2em}\=\kill
\textrm{\textbf{Function}} CollatzAEAD.Encrypt$(K, N, A, P)$:\\[4pt]
\textrm{\textit{Initialization:}}\\
\> 1. $s \leftarrow (0, \ldots, 0)$\\
\> 2. Load $N$ into $s[0..1]$; load $K$ into $s[4..7]$\\
\> 3. $s[7] \leftarrow s[7] \oplus \mathtt{DS\_AEAD}$\\
\> 4. $s \leftarrow \mathrm{CollatzPerm}(s, R_a)$\\
\> 5. $s[4..7] \leftarrow s[4..7] \oplus K$ \textrm{\small (post-init key injection)}\\[4pt]
\textrm{\textit{Process Associated Data $A$:}}\\
\> 6. \textbf{for each} padded $r$-bit block $A_i$ of $A$:\\
\> \> $s[0..3] \leftarrow s[0..3] \oplus A_i$\\
\> \> $s \leftarrow \mathrm{CollatzPerm}(s, R_b)$\\
\> 7. $s[7] \leftarrow s[7] \oplus \mathtt{DS\_AD\_END}$ \textrm{\small (domain separation)}\\[4pt]
\textrm{\textit{Encrypt Plaintext $P$:}}\\
\> 8. \textbf{for each} $r$-bit block $P_i$:\\
\> \> $C_i \leftarrow s[0..3] \oplus P_i$ \textrm{\small (keystream XOR)}\\
\> \> $s[0..3] \leftarrow C_i$ \textrm{\small (absorb ciphertext---duplex feedback)}\\
\> \> $s \leftarrow \mathrm{CollatzPerm}(s, R_b)$\\[4pt]
\textrm{\textit{Finalization:}}\\
\> 9. $s[4..7] \leftarrow s[4..7] \oplus K$\\
\> 10. $s \leftarrow \mathrm{CollatzPerm}(s, R_a)$\\
\> 11. $s[4..7] \leftarrow s[4..7] \oplus K$\\
\> 12. $T \leftarrow s[4..5]$ \textrm{\small (128-bit tag from capacity)}\\
\> 13. \textbf{return} $(C, T)$
\end{tabbing}
\end{quote}

Decryption follows the same structure with the inverse operation on the
plaintext blocks. The tag is verified using constant-time byte comparison
before releasing any plaintext.

\subsection{Security Properties}

\begin{itemize}
    \item \textbf{Confidentiality}: ciphertext is indistinguishable from
          random under CPA, bounded by $2^{128}$ queries.
    \item \textbf{Integrity}: forging a valid $(C, T)$ pair requires
          $2^{128}$ work (tag is extracted from the capacity).
    \item \textbf{Nonce misuse}: as with all nonce-based AEAD, reusing
          a nonce with the same key leaks information. This is a standard
          limitation shared by AES-GCM and ChaCha20-Poly1305.
\end{itemize}

% =====================================================================
\section{CollatzMAC: Message Authentication Code}
\label{sec:mac}
% =====================================================================

CollatzMAC is a keyed-sponge MAC \cite{sponge_security} with a 256-bit
key and configurable tag length (128 or 256 bits).

\begin{quote}\tt
\begin{tabbing}
\hspace{2em}\=\hspace{2em}\=\kill
\textrm{\textbf{Function}} CollatzMAC$(K, M, d)$:
\textrm{\small MAC with $d$-bit tag}\\[4pt]
\> 1. $s \leftarrow (0, \ldots, 0)$\\
\> 2. $s[7] \leftarrow s[7] \oplus \mathtt{DS\_MAC}$;\quad
$s[6] \leftarrow s[6] \oplus d$\\
\> 3. $s[0..3] \leftarrow s[0..3] \oplus K$ \textrm{\small (absorb key into rate)}\\
\> 4. $s \leftarrow \mathrm{CollatzPerm}(s, R_a)$\\
\> 5. \textbf{for each} padded $r$-bit block $M_i$ of $M$:\\
\> \> $s[0..3] \leftarrow s[0..3] \oplus M_i$\\
\> \> $s \leftarrow \mathrm{CollatzPerm}(s, R_b)$\\
\> 6. $s[4..7] \leftarrow s[4..7] \oplus K$ \textrm{\small (key injection before final)}\\
\> 7. $s \leftarrow \mathrm{CollatzPerm}(s, R_a)$\\
\> 8. Squeeze $d$ bits from rate\\
\> 9. \textbf{return} tag
\end{tabbing}
\end{quote}

Security is bounded by $\min(2^{c/2}, 2^k) = 2^{128}$ MAC queries,
following the keyed-sponge PRF bound \cite{sponge_security}.

% =====================================================================
\section{Domain Separation}
\label{sec:domain}
% =====================================================================

Each mode uses a distinct domain separation constant XORed into the
capacity before processing:

\begin{center}
\begin{tabular}{|l|l|}
\toprule
\textbf{Mode} & \textbf{Constant} \\
\midrule
Hash    & $\mathtt{0x01}$ \\
AEAD    & $\mathtt{0x02}$ \\
MAC     & $\mathtt{0x04}$ \\
AAD end & $\mathtt{0x80}$ \\
Final   & $\mathtt{0xFF}$ \\
\bottomrule
\end{tabular}
\end{center}

This ensures that hash digests, ciphertext/tags, and MAC tags computed
from the same input are cryptographically independent.

% =====================================================================
\section{Implementation}
\label{sec:impl}
% =====================================================================

\subsection{Reference Implementation}

The reference implementation is in C++17, structured as:
\begin{itemize}
    \item \texttt{collatz\_perm.h}: header-only permutation with
          architecture-specific paths.
    \item \texttt{collatz\_crypto.h}: public API for Hash, AEAD, MAC.
    \item \texttt{collatz\_crypto.cpp}: implementation, tests, and
          benchmarks.
\end{itemize}

\subsection{ARM64 Optimization}

On AArch64 (Apple Silicon), the \texttt{collatz\_mul\_pair} function is
implemented as 11 inline assembly instructions:

{\small
\begin{verbatim}
    eor  a, a, a, ror #41     // pre-mix: a ^= rotl(a, 23)
    orr  a, a, #1             // ensure odd
    add  t1, a, a, lsl #1     // t1 = 3*a
    add  t1, t1, #1           // t1 = 3*a + 1
    rbit t2, t1               // bit-reverse for ctz
    clz  t2, t2               // ctz = clz(rbit(x))
    lsr  a, t1, t2            // a = (3a+1) >> ctz
    mul  t1, a, b             // lo = a * b
    umulh t2, a, b            // hi = a * b
    eor  a, a, t1             // a ^= lo
    eor  b, b, t2             // b ^= hi
\end{verbatim}
}

The function uses only general-purpose registers and has no
data-dependent branches or memory accesses, making it constant-time with
respect to the state values.

\textbf{Note on \texttt{RBIT+CLZ}.} ARM64 does not have a native
count-trailing-zeros instruction. Instead, we compose bit-reverse
(\texttt{RBIT}) with count-leading-zeros (\texttt{CLZ}) to obtain
$\ctz$ in 2 cycles. On x86-64, the \texttt{TZCNT} instruction
provides this directly.

\subsection{Constant-Time Guarantees}

\begin{itemize}
    \item The Collatz step uses \texttt{ORR} (bitwise OR with 1) instead
          of a conditional branch.
    \item Tag verification uses byte-wise OR accumulation, not early
          return.
    \item No secret-dependent array indexing or memory access patterns.
\end{itemize}

\textbf{Caveat:} The variable-length right shift $t \gg \ctz(t)$
depends on the data-dependent shift amount $\ctz(t)$. On most modern
processors (ARM64, x86-64), barrel shifters execute variable shifts in
constant time. However, this should be verified on the target platform.

% =====================================================================
\section{Empirical Security Evaluation}
\label{sec:security}
% =====================================================================

We run 70 automated tests covering permutation diffusion, hash
avalanche, bit distribution, collision resistance, AEAD correctness and
integrity, and MAC correctness.

\subsection{Permutation Diffusion}
\label{sec:diffusion}

We verify that flipping any single bit in any of the 8 state words,
after 12 rounds of CollatzPerm-512, changes all 8 output words. We test
all 8 words at bit positions 0, 16, 32, and 48 (32 tests total).

\textbf{Result:} 32/32 tests pass (all 8 words change in every case).
This was achieved after adding the pre-mix rotation
$a \leftarrow a \oplus \rotl(a, 23)$ to prevent the $a \,|\, 1$
operation from absorbing bit 0.

\subsection{Avalanche Effect}

For each input string, we flip every bit of the input and measure the
Hamming distance between the original and modified hash digests,
expressed as a percentage of the output width:

\begin{center}
\begin{tabular}{|l|r|r|r|}
\toprule
\textbf{Input} & \textbf{256-bit} & \textbf{512-bit} & \textbf{1024-bit} \\
\midrule
\texttt{"hello world"}    & 50.3\% & 49.9\% & 50.0\% \\
\texttt{"sample input 12"} & 49.8\% & 50.1\% & 50.1\% \\
\texttt{"test input here"} & 49.6\% & 49.8\% & 49.9\% \\
\bottomrule
\end{tabular}
\end{center}

All values fall within the acceptable range of $50 \pm 5\%$, indicating
near-ideal avalanche behavior across all output widths. The ideal value
is 50\%, meaning each output bit flips with probability $\frac{1}{2}$
when any input bit changes.

\subsection{Bit Distribution}

For 50,000 sequential inputs (\texttt{"dist\_0"} through
\texttt{"dist\_49999"}), we measure the percentage of times each bit
position is set to 1:

\begin{center}
\begin{tabular}{|l|r|r|}
\toprule
\textbf{Output Width} & \textbf{Min} & \textbf{Max} \\
\midrule
256-bit  & 49.5\% & 50.6\% \\
1024-bit & 49.2\% & 50.6\% \\
\bottomrule
\end{tabular}
\end{center}

All bit positions are within $50 \pm 1\%$, indicating excellent
uniformity.

\subsection{Collision Resistance}

We hash $10^6$ distinct inputs at both 256-bit and 1024-bit output
widths and check for duplicate digests:

\begin{center}
\begin{tabular}{|l|r|r|}
\toprule
\textbf{Output Width} & \textbf{Inputs} & \textbf{Collisions} \\
\midrule
256-bit  & 1,000,000 & 0 \\
1024-bit & 1,000,000 & 0 \\
\bottomrule
\end{tabular}
\end{center}

\textbf{Note:} A birthday-bound collision search for a 256-bit hash
requires approximately $2^{128}$ inputs; our test of
$10^6 \approx 2^{20}$ inputs is insufficient to find collisions by
chance and serves only to verify the absence of trivial implementation
errors.

\subsection{AEAD Integrity}

We test authentication integrity by verifying that decryption fails
(returns $\bot$) when any of the following are modified:

\begin{center}
\begin{tabular}{|l|r|}
\toprule
\textbf{Modification} & \textbf{Result} \\
\midrule
Flip 1 ciphertext bit & Rejected \\
Flip 1 tag bit         & Rejected \\
Flip 1 AAD bit         & Rejected \\
Wrong key (1-bit diff) & Rejected \\
Wrong nonce (1-bit diff) & Rejected \\
\bottomrule
\end{tabular}
\end{center}

All 5 authentication tests pass. AEAD roundtrip correctness is verified
for plaintext lengths 0, 1, 15, 16, 31, 32, 33, 63, 64, 100, and 1000
bytes (all 11 tests pass).

\subsection{Key Sensitivity}

Two 256-bit keys differing by a single bit produce ciphertexts with 93\%
of bytes differing and completely different authentication tags.

% =====================================================================
\section{Performance}
\label{sec:perf}
% =====================================================================

All benchmarks were run on Apple M3 Max (ARM64, 10 performance cores)
with \texttt{g++ -O3 -march=native}. All comparison implementations are
standalone C++ (no hardware acceleration extensions, no external
libraries).

\subsection{Hash Performance}

\begin{center}
\begin{tabular}{|l|l|r|r|}
\toprule
\textbf{Input} & \textbf{Algorithm} & \textbf{ops/sec} & \textbf{ns/op} \\
\midrule
5 B  & CollatzHash-256  & 17,066,218 & 58.6 \\
     & CollatzHash-1024 & 5,963,157  & 167.7 \\
     & SHA-256          & 4,828,820  & 207.1 \\
\midrule
12 B & CollatzHash-256  & 16,745,487 & 59.7 \\
     & CollatzHash-1024 & 5,956,764  & 167.9 \\
     & SHA-256          & 4,962,137  & 201.5 \\
\midrule
43 B & CollatzHash-256  & 10,903,183 & 91.7 \\
     & CollatzHash-1024 & 4,944,270  & 202.3 \\
     & SHA-256          & 4,911,992  & 203.6 \\
\bottomrule
\end{tabular}
\end{center}

CollatzHash-256 is \textbf{3.4$\times$ faster} than SHA-256 on short
inputs. The advantage narrows for longer inputs as the sponge
absorption cost grows proportionally.

\subsection{AEAD Performance}

\begin{center}
\begin{tabular}{|r|r|r|l|}
\toprule
\textbf{Size} & \textbf{CollatzAEAD} & \textbf{ChaCha20} & \textbf{Ratio} \\
\midrule
64 B   & 332 MB/s   & 917 MB/s  & 0.36$\times$ \\
256 B  & 629 MB/s   & 951 MB/s  & 0.66$\times$ \\
1024 B & 800 MB/s   & 931 MB/s  & 0.86$\times$ \\
4096 B & 859 MB/s   & 942 MB/s  & 0.91$\times$ \\
\bottomrule
\end{tabular}
\end{center}

\textbf{Note:} ChaCha20 here is encryption-only (no Poly1305
authentication). CollatzAEAD provides combined encrypt+authenticate in a
single pass, so the true comparison against ChaCha20-Poly1305 would be
more favorable.

\subsection{MAC Performance}

\begin{center}
\begin{tabular}{|l|r|r|}
\toprule
\textbf{Algorithm} & \textbf{ops/sec} & \textbf{ns/op} \\
\midrule
CollatzMAC-128 & 6,941,647  & 144.1 \\
HMAC-SHA-256   & 1,290,474  & 774.9 \\
\bottomrule
\end{tabular}
\end{center}

CollatzMAC is \textbf{5.4$\times$ faster} than HMAC-SHA-256.

% =====================================================================
\section{Comparison with Existing Primitives}
\label{sec:comparison}
% =====================================================================

\begin{center}
\begin{tabular}{|l|c|c|c|c|}
\toprule
\textbf{Property} & \textbf{SHA-256} & \textbf{SHA-3-256} & \textbf{Ascon} & \textbf{Collatz} \\
\midrule
State width        & 256 b   & 1600 b  & 320 b  & 512 b \\
Capacity           & ---     & 512 b   & 192 b  & 256 b \\
Security target    & 128 b   & 128 b   & 128 b  & 128 b \\
Nonlinearity       & Ch, Maj & $\chi$  & S-box  & Collatz+mul \\
Rounds/block       & 64      & 24      & 8      & 8 \\
Output widths      & 256     & XOF     & ---    & 256/512/1024 \\
AEAD mode          & No      & No      & Yes    & Yes \\
MAC mode           & HMAC    & KMAC    & Yes    & Yes \\
Standardized       & NIST    & NIST    & NIST   & No \\
Peer review        & 20+ yr  & 15+ yr  & 10+ yr & This paper \\
\bottomrule
\end{tabular}
\end{center}

% =====================================================================
\section{Analysis of the Collatz Nonlinearity}
\label{sec:analysis}
% =====================================================================

\subsection{Why the Collatz Step Provides Nonlinearity}

Traditional symmetric ciphers derive nonlinearity from S-boxes (lookup
tables with high differential and linear branch numbers) or from
algebraic operations like modular multiplication. The Collatz step
$f(a) = (3a+1) \gg \ctz(3a+1)$ provides nonlinearity through:

\begin{enumerate}
    \item \textbf{Data-dependent shift.} The shift amount
          $\ctz(3a+1)$ varies unpredictably with $a$. This means the
          output bits are a data-dependent permutation of the input bits
          combined with the multiplication by 3, providing high algebraic
          degree.
    \item \textbf{Bit truncation.} The right shift discards
          $\ctz(3a+1)$ low-order bits, making the function non-invertible
          on the full 64-bit word. This loss of information is analogous
          to the compression in S-box designs.
    \item \textbf{Interaction with multiplication.} The subsequent
          128-bit multiply amplifies the nonlinear bit patterns across
          both output words.
\end{enumerate}

\subsection{Differential Properties}

Let $\Delta_a$ denote an input difference in word $a$. After the
Collatz step:
$$f(a) = (3(a \,|\, 1) + 1) \gg \ctz(3(a \,|\, 1) + 1)$$

The output difference $f(a \oplus \Delta_a) \oplus f(a)$ depends on
both the algebraic carry propagation of the $3a+1$ operation and the
data-dependent shift. A formal differential analysis would enumerate
the probability $\Pr[\Delta_{\mathrm{out}} | \Delta_{\mathrm{in}}]$
over all possible inputs, which we leave as future work.

\subsection{The Pre-Mix Rotation}

Without the pre-mix rotation $a \leftarrow a \oplus \rotl(a, 23)$,
the $a \,|\, 1$ operation absorbs (destroys) bit 0 of the input. This
creates a specific weakness: for any state word, toggling bit 0 has no
effect after the Collatz step, violating the strict avalanche criterion.

The pre-mix XORs bit 0 into bit 23, bit 1 into bit 24, etc., ensuring
that every input bit---including bit 0---influences the Collatz step
through a different bit position. The rotation distance 23 (equivalently,
\texttt{ROR \#41} on ARM64) is chosen as:
\begin{itemize}
    \item Coprime to 64 (avoiding periodic alignment),
    \item Not equal to any of the linear layer rotation constants
          (19, 31, 7, 53).
\end{itemize}

% =====================================================================
\section{Discussion}
\label{sec:discussion}
% =====================================================================

\subsection{Strengths}

\begin{enumerate}
    \item \textbf{Novel nonlinearity.} The Collatz step introduces a
          number-theoretic nonlinear operation not found in any existing
          cryptographic primitive. The chaotic behavior of the $3n+1$
          mapping may provide resistance to algebraic attacks, though
          this requires further analysis.
    \item \textbf{Speed.} The Collatz--multiply hybrid achieves better
          throughput than SHA-256 for hashing and is competitive with
          ChaCha20 for encryption, without requiring hardware
          acceleration.
    \item \textbf{Simplicity.} One permutation serves all three modes.
          The implementation is under 500 lines of C++ with no external
          dependencies.
    \item \textbf{Standard constructions.} The sponge and duplex modes
          inherit formal security proofs from the literature, reducing
          the analysis burden to the permutation itself.
\end{enumerate}

\subsection{Limitations and Open Questions}

\begin{enumerate}
    \item \textbf{No formal cryptanalysis.} The security of the
          permutation rests on empirical evidence. Formal analysis of
          differential, linear, and algebraic properties is needed.
    \item \textbf{Algebraic structure.} The Collatz step is
          algebraically simple ($3n+1$ followed by right-shift). An
          attacker with knowledge of the algebraic structure might find
          exploitable properties. The 128-bit multiply mitigates this but
          does not eliminate the concern.
    \item \textbf{Bit-0 absorption.} The $a \,|\, 1$ operation destroys
          one bit of input per Collatz step. The pre-mix rotation
          recovers this information, but the interaction between the
          rotation and the multiply deserves further study.
    \item \textbf{No hardware acceleration.} Unlike AES (AES-NI) and
          SHA-256 (SHA extensions), there is no dedicated hardware
          support for Collatz operations.
    \item \textbf{Not standardized.} Adoption in production systems
          requires peer review and potential standardization.
    \item \textbf{Variable shift timing.} While barrel shifters on
          modern processors execute in constant time regardless of shift
          amount, some microcontrollers may not. Deployments on embedded
          targets should verify constant-time behavior.
\end{enumerate}

\subsection{Recommended Use Cases}

\begin{itemize}
    \item File integrity checking and deduplication.
    \item API authentication tokens and session management.
    \item Encrypted storage and communication in applications where
          SHA-256 is the performance bottleneck and hardware AES is
          unavailable.
    \item Research into number-theoretic cryptographic primitives.
\end{itemize}

\subsection{Not Recommended (Without Further Analysis)}

\begin{itemize}
    \item Digital signatures for high-value financial transactions.
    \item Government or military classified communications.
    \item Any application requiring FIPS~140 or Common Criteria
          certification.
\end{itemize}

% =====================================================================
\section{Future Work}
\label{sec:future}
% =====================================================================

\begin{enumerate}
    \item \textbf{Formal differential/linear cryptanalysis} of
          CollatzPerm-512, including computation of the best differential
          and linear characteristics.
    \item \textbf{Algebraic attack resistance}: analyze the algebraic
          degree of the round function and resistance to Gr\"obner basis
          attacks.
    \item \textbf{Bounds on round count}: determine the minimum number
          of rounds for full diffusion and provable resistance to known
          attacks.
    \item \textbf{x86-64 SIMD optimization}: exploit AVX2/AVX-512 for
          parallel Collatz-multiply pairs.
    \item \textbf{Nonce-misuse resistance}: design an SIV-like mode
          (Synthetic IV) that provides security even under nonce reuse.
    \item \textbf{Test vectors}: publish a comprehensive set of
          known-answer test vectors for interoperability.
    \item \textbf{Side-channel analysis}: formal verification of
          constant-time behavior on various microarchitectures.
\end{enumerate}

% =====================================================================
\section{Conclusion}
\label{sec:conclusion}
% =====================================================================

We have presented a complete cryptographic suite built on the Collatz
mapping: a 512-bit permutation, a sponge hash, a duplex-sponge AEAD
cipher, and a keyed-sponge MAC. The suite targets 128-bit security,
outperforms SHA-256 by 3.4$\times$ for hashing, matches ChaCha20 for
large-message encryption, and surpasses HMAC-SHA-256 by 5.4$\times$ for
message authentication.

The design demonstrates that the chaotic behavior of the $3n+1$ mapping
can serve as an effective nonlinear mixing primitive when combined with
128-bit integer multiplication for diffusion. We invite the cryptographic
community to analyze the permutation's resistance to differential,
linear, and algebraic attacks.

The reference implementation is available at:
\texttt{https://github.com/zmockwebdesign/collatz-v2}

% =====================================================================
\section*{Acknowledgments}
% =====================================================================

The sponge and duplex constructions follow the frameworks established by
Bertoni, Daemen, Peeters, and Van Assche for Keccak/SHA-3 and by
Dobraunig, Eichlseder, Mendel, and Schl\"affer for Ascon.

% =====================================================================
\begin{thebibliography}{99}

\bibitem{lagarias1985}
J.~C. Lagarias,
``The $3x+1$ problem and its generalizations,''
\textit{American Mathematical Monthly}, vol.~92, no.~1, pp.~3--23, 1985.

\bibitem{keccak_sponge}
G.~Bertoni, J.~Daemen, M.~Peeters, and G.~Van Assche,
``Cryptographic sponge functions,''
\texttt{https://keccak.team/files/CSF-0.1.pdf}, 2011.

\bibitem{sha3}
NIST,
``SHA-3 Standard: Permutation-Based Hash and Extendable-Output
Functions,'' FIPS~202, 2015.

\bibitem{ascon}
C.~Dobraunig, M.~Eichlseder, F.~Mendel, and M.~Schl\"affer,
``Ascon v1.2: Lightweight Authenticated Encryption and Hashing,''
\textit{Journal of Cryptology}, vol.~34, 2021.
NIST Lightweight Cryptography Standard, 2023.

\bibitem{sponge_security}
G.~Bertoni, J.~Daemen, M.~Peeters, and G.~Van Assche,
``On the security of the keyed sponge construction,''
\textit{SKEW 2011}, 2011.

\bibitem{duplex}
G.~Bertoni, J.~Daemen, M.~Peeters, and G.~Van Assche,
``Duplexing the sponge: Single-pass authenticated encryption and other
applications,'' \textit{SAC 2011}, LNCS~7118, pp.~320--337, 2012.

\bibitem{rapidhash}
N.~Kurz,
``rapidhash --- Very fast, high quality, platform-independent hashing
algorithm,'' \texttt{https://github.com/Nicoshev/rapidhash}, 2024.

\bibitem{wyhash}
Y.~Wang,
``wyhash --- The fastest hash function,''
\texttt{https://github.com/wangyi-fudan/wyhash}, 2023.

\bibitem{xxhash}
Y.~Collet,
``xxHash --- Extremely fast non-cryptographic hash algorithm,''
\texttt{https://github.com/Cyan4973/xxHash}, 2024.

\end{thebibliography}

% =====================================================================
\appendix
\section{Test Vectors}
\label{app:testvectors}

The following test vectors are produced by the reference implementation.

\subsection{CollatzHash-256}

\begin{verbatim}
Input:  "" (empty)
Output: f266161215a29b9f62dd4b71b4b05144
        40841de8358cd66db10fe489202ccf58

Input:  "abc"
Output: a170ec42f3e1513522b7513f478272db
        054729f09bcd9e4d30f6a661c7a8e0cd

Input:  "hello world"
Output: 4bec83a5598cecee53dc89697f51e928
        e93783237399d7b03415bb6c0ed08c9c
\end{verbatim}

\subsection{CollatzHash-512}

\begin{verbatim}
Input:  "" (empty)
Output: f266161215a29b9f62dd4b71b4b05144
        40841de8358cd66db10fe489202ccf58
        1a14744ac45936e49b4e751374281ed7
        72b5b1e6be4811fe31015171e22c5812

Input:  "abc"
Output: a170ec42f3e1513522b7513f478272db
        054729f09bcd9e4d30f6a661c7a8e0cd
        13d531694946463d91e759c1acc0d320
        ad5c3995e8a9b32a1b5351dfd95ffe79

Input:  "hello world"
Output: 4bec83a5598cecee53dc89697f51e928
        e93783237399d7b03415bb6c0ed08c9c
        36d90399456178f7a3a6082e9620af31
        e11bd7c81b179784fd351a9ba58fa378
\end{verbatim}

\subsection{CollatzMAC-128}

\begin{verbatim}
Key:    00000000000000000000000000000000
        00000000000000000000000000000000

Input:  "" (empty)
Output: daff09bb11ac65519bfbe22483fed036

Input:  "abc"
Output: 951eebfd00e05e79c5d182e11b2a3286
\end{verbatim}

\end{document}
